\section{Conclusion}
\label{sec:conclusion}

We presented \spork{}.
Our central finding is that agentic LLM workloads exhibit exploitable early intent: the
running model already knows its next tool call before finishing its reasoning, with
fork-at-start name accuracy of 74.6--99.6\% on Qwen3-32B against a tool wait of
16--37\% of agent wall time.
\spork{} turns this into latency with the lightest possible machinery, requiring no
draft model, no historical traces, and no retraining: D1--D3 map the early-intent
signals onto EQ1's terms, and our experiments validate the cost model on real web tools
(GAIA), short-latency APIs (HotpotQA), and the controlled tau2 latency sweep.
The mechanism generalizes across scale and architecture, from Qwen3-4B to Qwen3-32B and
a 3B-active mixture-of-experts, while preserving task accuracy.

\parabf{Limitations.}
\spork{} speculates only read-only tools today; supporting write operations is future
work, requiring sandboxing or agent-state checkpointing so a mispredicted speculative
call can be rolled back.
It relies on open-weight models served by an engine that exposes per-token logprobs and
a completion endpoint; closed-source APIs cannot host the fork thread.
Finally, we target latency-oriented serving with spare capacity; under heavy batching,
probe traffic competes with foreground decode and the overlap benefit shrinks.

\parabf{Future work.}
Sandboxed or checkpointed execution would extend speculation to write tools: a
mispredicted call rolls back to the pre-speculation agent state.
Deeper engine integration could schedule probes into idle batch slots, shrinking
$\Toh$ further and widening the favorable regime.
